\section{Two distinct finite minimal covers}
Use $a,b,c$ for rotation generators, with common relations
\begin{equation}\label{eq:commonrels}
                         (ab)^2=(bc)^2=(abc)^2=1.
\end{equation}
Let $P$ have, in addition, the literal relations
\begin{equation}\label{eq:Prels}
 a^6=b^9=c^6=1,\qquad ba^2=a^2b^4,\qquad c^2b=b^4c^2.
\end{equation}
Let $T$ have \eqref{eq:commonrels} and
\begin{equation}\label{eq:Trels}
 a^6=b^9=c^2=1,\qquad (ab^{-2})^2=1,\qquad
                         ab^2a(ab^{-1}a)^2=1,
\end{equation}
and let $T^\delta$ have \eqref{eq:commonrels} and
\begin{equation}\label{eq:Tdrels}
 a^2=b^9=c^6=1,\qquad(c^{-1}b^2)^2=1,\qquad
                         bc^2b(c^{-1}bc^{-1})^2=1.
\end{equation}
These are the three presentations printed on page 7 of the checked author edition. No extra relation is added to make an order match. We set
\[
 R=P\diamond\overline P,
       \qquad Q=R\diamond T,
       \qquad Q'=R\diamond T^\delta,
\]
where the mirror marking of $P$ is $(a^{-1},a^2b,c)$.

The complete required program \texttt{chiral\_cover.g} starts with the three finitely presented groups, obtains exact GAP isomorphisms to finite permutation groups, and generates the mixes inside direct products from the literal paired generators. It verifies the following table. Every order in the table is a rotation-group order; ``mirror'' means that the displayed mirror assignment extends to an automorphism.
\begin{center}\small
\begin{tabular}{@{}lrrrrl@{}}\toprule
Group & Order & Orders of $a,b,c$ & End orders & End intersection & Mirror\\\midrule
$P$ &324&$6,9,6$&$54,54$&9&no\\
$T$ &324&$6,9,2$&$162,18$&9&yes\\
$T^\delta$ &324&$2,9,6$&$18,162$&9&yes\\
$R$ &972&$6,9,6$&$162,162$&27&yes\\
$Q$ &2916&$6,9,6$&$162,162$&9&yes\\
$Q'$ &2916&$6,9,6$&$162,162$&9&yes\\\bottomrule
\end{tabular}
\end{center}
For every row the program also checks the rotation relations, $|\langle b\rangle|=9$, and
\[
       \langle a\rangle\cap\langle b\rangle=1,
             \qquad\langle b\rangle\cap\langle c\rangle=1.
\]
For the five successful rows it checks the actual subgroup equality
$\langle a,b\rangle\cap\langle b,c\rangle=\langle b\rangle$; for $R$ it checks failure. These equalities are not inferred from similarly sized but unidentified groups. The table therefore proves, using the retained recursive correspondence, that $P$ is chiral, $T,T^\delta,Q,Q'$ are directly regular polytopes, and $R$ is not a polytope. For $P$, the specified mirror assignment fails to define a homomorphism. For every positive row the assignment is checked to be bijective.

The program computes the first-coordinate images of the second-projection kernels in the two ends of $R$. Both are exactly $\langle b^3\rangle\leq\Gamma(P)$, of order three. Thus the finite example exhibits directly the obstruction in Theorem~\ref{thm:covercriterion}: its medial polygon has trivial chirality, whereas $b^3$ lies nontrivially in both end chirality groups.

It also verifies a marked isomorphism from the facet of $T$ to the facet of $R$, and the dual isomorphism
\[
                (a,b,c)_T\longmapsto(c^{-1},b^{-1},a^{-1})_{T^\delta}.
\]
The facet subgroup of $T$ has index two, so the directly regular polytope $T$ has exactly two facets. Under the standard interpretation of the trivial extension used by the source, this is the trivial extension of that facet. The dual map proves the claimed identity of $T^\delta$, not just an agreement of orders.

The marked epimorphisms
\[
                         R\to P,\qquad Q\to R,
                         \qquad Q'\to R
\]
are verified and each has kernel order three. The common invariant-kernel theorem proves smallestness of $R$; the common prime-degree lemma then proves minimality of $Q$ and $Q'$ as regular polytope covers of $P$. Doubling rotation orders gives the source's respective flag counts $648$, $1944$ and $5832$.

To distinguish the two minimal covers, the program supplies the word
\begin{equation}\label{eq:distinguish}
                            w=(bc^{-1})^2ac^2a^{-1}.
\end{equation}
Its order is $1$ in $Q$ and $3$ in $Q'$. Hence no marked isomorphism exists. This also excludes an unbased rank-preserving polytope isomorphism. Such an isomorphism commutes with every flag-adjacency operation. Composing it with an automorphism of the regular target carries the image of a chosen base flag back to that target's chosen base flag; it would then identify all distinguished rotation generators. It would preserve the value and order of \eqref{eq:distinguish}, a contradiction. No assertion about dual isomorphisms follows from this word test.

The calculation retains GAP's exact algorithms for finitely presented groups, permutation isomorphisms, subgroup intersections and homomorphisms. It uses neither RAMP nor a group catalogue. The source used RAMP for its example; the replacement program is the previously reviewed bare-presentation computation, copied in full. Its inputs and logical outputs are all specified above. The uniform criterion and the two infinite-family arguments do not depend on the finite table.
